documentclass[a4paper,12pt]{article}

% Настройка русского языка и шрифтов
\usepackage[T2A]{fontenc}
\usepackage[utf8]{inputenc}
\usepackage[russian]{babel}

% Библиотеки для красивой математики и нумерации формул
\usepackage{amsmath}
\usepackage{amsfonts}
\usepackage{amssymb}

% Настройка полей страницы
\usepackage[left=2cm,right=2cm,top=2cm,bottom=2cm]{geometry}

\begin{document}

\pagestyle{empty}

\begin{center}
    \large\bfseries Численное моделирование одномерного уравнения переноса
\end{center}

\vspace{0.5cm}

В данной работе рассматривается программная реализация трех разностных схем для решения линейного уравнения переноса первого порядка (уравнения конвекции):
\begin{equation*}
    \frac{\partial u}{\partial t} + c \frac{\partial u}{\partial x} = 0
\end{equation*}

В качестве тестового начального условия $u(0, x) = f(x)$ используется Гауссова функция. Эта задача является линейной и имеет решение $u(x, t) = g(x - ct), t > 0$.

\section{Реализованные разностные схемы}

Ниже приведены математические аппроксимации схем, взятые из учебника \textit{В. Зализняка «Основы вычислительной физики. Часть 1. Введение в конечно-разностные методы» Глава 3}:

\subsection{Схема 3.3}
Схемы устойчивы только при выполнении условия Куранта-Фридрикса-Леви: $K = \frac{c \cdot \tau}{h} \le 1$.
\begin{itemize}
    \item[] {Разностное уравнение (в исходном виде):}
    \begin{equation*}
        \frac{u_{n}^{k+1} - u_{n}^{k}}{\tau} + c \frac{u_{n}^{k} - u_{n-1}^{k}}{h} = 0
    \end{equation*}
    
    \item[] {Расчетная формула в коде (\texttt{dif\_scheme\_1}):}
    \begin{equation*}
        u_{n}^{k+1} = u_{n}^{k} - K \cdot (u_{n}^{k} - u_{n-1}^{k})
    \end{equation*}
\end{itemize}

\subsection{Схема 3.5}

\begin{itemize}
    \item[] {Разностное уравнение (в исходном виде):}
    \begin{equation*}
        \frac{u_{n}^{k+1} - u_{n}^{k}}{\tau} + c \frac{u_{n}^{k+1} - u_{n-1}^{k+1}}{h} = 0
    \end{equation*}
    
    \item[] {Расчетная формула в коде (\texttt{dif\_scheme\_2}):}
    \begin{equation*}
        u_{n}^{k+1} = \frac{u_{n}^{k} + K \cdot u_{n-1}^{k+1}}{1 + K}
    \end{equation*}
\end{itemize}

\subsection{Схема 3.6}
Поскольку при $c > 0$ её шаблон направлен против потока, схема является  неустойчивой.

\begin{itemize}
    \item[] {Разностное уравнение (в исходном виде):}
    \begin{equation*}
        \frac{u_{n}^{k+1} - u_{n}^{k}}{\tau} - c \frac{u_{n+1}^{k} - u_{n}^{k}}{h} = 0
    \end{equation*}
    
    \item[] {Расчетная формула в коде (\texttt{dif\_scheme\_3}):}
    \begin{equation*}
        u_{n}^{k+1} = u_{n}^{k} - K \cdot (u_{n+1}^{k} - u_{n}^{k})
    \end{equation*}
\end{itemize}

\end{document}
